\section*{Executive Summary}
\addcontentsline{toc}{section}{Executive Summary}

This research proposal presents the development of the Multimodal and Context-Aware Ensemble
(MACE) framework for the detection of inauthentic social media content. The proposal is based on
the increasing impact of misinformation on public trust, online information credibility, and
decision-making. As misinformation is now commonly shared through text, images, reposts,
comments, user interactions, and propagation patterns, the detection of inauthentic content
requires more than text-only analysis.

The main purpose of this research is the investigation of whether the replacement of a text-only
content branch in existing context-aware models, specifically DANES and GETAE, with a multimodal
content branch can improve misinformation detection performance. The proposed MACE framework keeps
the original social-context branch from these models while replacing the text-only branch with a
multimodal branch based on HMCAN, enabling the processing of both textual and visual information
while maintaining the existing propagation-based context structure.

The research is guided by three main questions: whether multimodal substitution improves detection
performance and at what computational cost; the relative contribution of the multimodal content
branch versus the social-context branch through ablation testing; and the dataset adaptation
requirements for using existing multimodal social media datasets within a DANES/GETAE-style
ensemble structure.

The proposed methodology follows a staged experimental design. Stage 1 replicates DANES and GETAE
on Twitter15 and Twitter16. Stage 2 assesses candidate datasets such as FakeNewsNet and MuMin for
text availability, image availability, label compatibility, graph suitability, class balance, and
computational feasibility. Stage 3 conducts controlled evaluation across multiple experimental
conditions, including text-only baselines, multimodal-only ablations, and the full MACE framework.

Evaluation uses standard classification metrics (accuracy, precision, recall, and F1-score)
alongside computational efficiency metrics (inference time, training time, peak memory usage,
storage requirements, and hardware configuration), assessing not only whether MACE improves
performance but whether the improvement justifies the additional computational cost.
